\begin{frame}{자주 하는 실수 1: ``$h_t$는 $x_t$만 의존한다''}
  RNN을 처음 접할 때 가장 흔한 오해: 순전파 수식에서
  $W_{hh}\, h_{t-1}$ 항을 ``보조적인 정규화 항''쯤으로 여기는 것.
  \vskip0.6em
  \begin{itemize}
    \item $h_t$는 \textbf{과거 전체 $(x_1, \ldots, x_t)$의 함수}이며,
      $x_t$ 하나만 의존하지 않는다
    \item ``은닉 상태가 과거를 기억한다''는 이 절의 모든 주장이
      \textbf{이 한 항}에 달려 있다
    \item $W_{hh} = 0$으로 두면:
      $h_t = \tanh(W_{xh}\, x_t + b_h)$로
      \textbf{``시점을 독립적으로 처리하는 MLP''로 퇴화}
    \item 결과: 10단어 문장의 10번째 단어를 예측할 때 \textbf{앞의
      9단어를 완전히 잊는} 모델
  \end{itemize}
  \vskip0.6em
  \kq{$W_{hh}\, h_{t-1}$을 빼면 RNN은 RNN이 아니다.}
\end{frame}
